Put \(d_0=p+2p^2\), \(N_n=n-2\), and \(U_k=Y_k-\xi\).  Constants below may depend on \(p,\Delta,\kappa\), but not on \(\theta\in\mathcal M_{\Delta,p}(\kappa)\), on time indices, or on coordinate indices.  The strengthened \(\text{lem:uniform-bin-eighth-moment}\) gives
\[
 \sup_\theta\mathbb E_\theta\lVert X_0\rVert_2^{16}<\infty.
\tag{1}
\]
Every coordinate of \(Y_0\) is either a coordinate of \(X_0\) or a product of one coordinate of \(X_0\) and one coordinate of \(X_1\) or \(X_2\).  Cauchy--Schwarz and stationarity therefore turn (1) into
\[
 \sup_\theta\mathbb E_\theta\lVert Y_0\rVert_2^8<\infty,
 \qquad
 \sup_\theta\mathbb E_\theta\lVert U_0\rVert_2^8<\infty.
\tag{2}
\]
The fixed-law versions follow from the all-orders part of that lemma.  By \(\text{lem:finite-depth-geometric-mixing}\), after increasing a constant to cover lags zero, one may fix \(\rho=e^{-c_{\mathrm{mix}}\Delta}\in(0,1)\) such that
\[
 \alpha_{Y,\theta}(h)\leq C_0\rho^h,
 \qquad h\geq0,
\tag{3}
\]
uniformly on the margin class.  A fixed law has the analogous bound with law-dependent \(C_0\) and \(\rho\).

We first prove the product inequality needed for every covariance and fourth-cumulant estimate below.  Let \(V\) and \(W\) be measurable with respect to two sigma-fields whose strong-mixing coefficient is at most \(a\), and suppose \(V\in L^s\), \(W\in L^t\), where \(s,t>1\) and \(s^{-1}+t^{-1}<1\).  Then
\[
 |\operatorname{Cov}(V,W)|
 \leq C_{s,t}\lVert V\rVert_s\lVert W\rVert_t
 a^{,1-s^{-1}-t^{-1}}.
\tag{4}
\]
To derive (4), first normalize both norms to one and split each variable into its positive and negative parts.  For nonnegative \(V,W\), layer cake and Fubini give
\[
 \operatorname{Cov}(V,W)=
 \int_0^\infty\!\int_0^\infty
 \{P(V>x,W>y)-P(V>x)P(W>y)\}\,dx\,dy.
\tag{5}
\]
The absolute value of the integrand is bounded simultaneously by \(a\), \(2P(V>x)\), and \(2P(W>y)\).  Markov's inequality consequently bounds the absolute value of (5) by
\[
 \int_0^\infty\!\int_0^\infty
 \min\{a,2x^{-s},2y^{-t}\}\,dx\,dy.
\tag{6}
\]
For \(a>0\), set \(A=(2/a)^{1/s}\) and \(B=(2/a)^{1/t}\).  On \([0,A]\times[0,B]\), (6) is at most \(aAB\).  On \((A,\infty)\times[0,B]\) it is at most \(2B\int_A^\infty x^{-s}dx\), and the symmetric rectangle has the analogous bound.  Each is a constant times \(a^{1-s^{-1}-t^{-1}}\).  On the last rectangle, the changes of variables \(x=Au\), \(y=Bv\) leave this same power of \(a\) times
\[
 \int_1^\infty\!\int_1^\infty\min\{u^{-s},v^{-t}\}\,du\,dv<\infty;
\tag{7}
\]
splitting at \(v=u^{s/t}\) proves finiteness precisely because \(s^{-1}+t^{-1}<1\).  Summing the four sign combinations proves (4).  If \(a=0\), the integrand in (5) vanishes and (4) follows by continuity.  This is an elementary proof of (4), not an additional mixing assumption.

Taking \(s=t=4\) in (4), and using (2)--(3), gives for every pair of scalar coordinates
\[
 |\operatorname{Cov}_\theta(U_0^a,U_h^b)|
 \leq C\rho^{|h|/2}.
\tag{8}
\]
Thus all coordinate covariance sequences are absolutely summable, uniformly.  Since \(d_0\) is fixed,
\[
 \sup_\theta\sum_{h\in\mathbb Z}
 \lVert\operatorname{Cov}_\theta(Y_0,Y_h)\rVert_{\mathrm{op}}<\infty.
\tag{9}
\]
In particular \(\Omega(\theta)\) exists and is uniformly finite.

We next supply the four-point calculation.  Fix ordered integers \(s_1\leq s_2\leq s_3\leq s_4\), scalar coordinates \(a_1,\ldots,a_4\), and put \(V_\ell=U_{s_\ell}^{a_\ell}\).  Let
\[
 G=\max_{1\leq j\leq3}(s_{j+1}-s_j),
\]
and choose \(j\) attaining the maximum.  For the left and right products, generalized H\"older and (2) give
\[
 \left\lVert\prod_{\ell=1}^{j}V_\ell\right\rVert_{8/j}\leq C,
 \qquad
 \left\lVert\prod_{\ell=j+1}^{4}V_\ell\right\rVert_{8/(4-j)}\leq C.
\tag{10}
\]
The two reciprocal exponents in (10) sum to \(1/2\).  Applying (4) across the gap \(G\) therefore gives
\[
 \left|\operatorname{Cov}\left(
 \prod_{\ell=1}^{j}V_\ell,
 \prod_{\ell=j+1}^{4}V_\ell\right)\right|
 \leq C\rho^{G/2}.
\tag{11}
\]
Here and below, coincident times, for which \(G=0\), are covered directly by (2).

All \(V_\ell\) are centered, so
\[
 \operatorname{cum}(V_1,V_2,V_3,V_4)
 =\mathbb E(V_1V_2V_3V_4)
 -\gamma_{12}\gamma_{34}-\gamma_{13}\gamma_{24}-\gamma_{14}\gamma_{23},
 \qquad \gamma_{ab}=\mathbb E(V_aV_b).
\tag{12}
\]
If \(j=1\), the first term in (12) is \(\operatorname{Cov}(V_1,V_2V_3V_4)\), which is controlled by (11); in each of the three pairing terms, the covariance containing \(V_1\) crosses the gap \(G=s_2-s_1\), so (8) controls that factor and (2) bounds the other.  If \(j=3\), the same cancellation holds with \(V_4\) separated from \(V_1V_2V_3\).  If \(j=2\), subtracting \(\gamma_{12}\gamma_{34}\) turns the first two terms of (12) into \(\operatorname{Cov}(V_1V_2,V_3V_4)\), controlled by (11), while every covariance in each of \(\gamma_{13}\gamma_{24}\) and \(\gamma_{14}\gamma_{23}\) crosses \(s_3-s_2=G\).  Thus the \(1|3\), \(2|2\), and \(3|1\) cases all yield the uniform envelope
\[
 |\operatorname{cum}(U_{s_1}^{a_1},U_{s_2}^{a_2},U_{s_3}^{a_3},U_{s_4}^{a_4})|
 \leq C\rho^{G/2}.
\tag{13}
\]

Let \(I\) be an integer interval of length \(\ell\), let \(a\in\mathbb R^{d_0}\) satisfy \(\lVert a\rVert_2=1\), and set \(S_I=\sum_{k\in I}a^{\mathsf T}U_k\).  Equation (9) gives \(\operatorname{Var}(S_I)\leq C\ell\).  By multilinearity of cumulants and (13), ordering four indices and writing their three adjacent gaps as \(d_1,d_2,d_3\) gives
\[
 |\operatorname{cum}_4(S_I)|
 \leq 4!C\ell
 \sum_{d_1,d_2,d_3\geq0}\rho^{\max(d_1,d_2,d_3)/2}
 \leq C'\ell.
\tag{14}
\]
The series is finite because the number of triples with maximum \(q\) is \((q+1)^3-q^3=O((q+1)^2)\).  The centered fourth-moment identity, obtained by expanding the fourth derivative of the cumulant generating function, is \(\mathbb E S_I^4=\operatorname{cum}_4(S_I)+3\operatorname{Var}(S_I)^2\).  Hence
\[
 \sup_\theta\mathbb E_\theta
 \left|\sum_{k\in I}a^{\mathsf T}U_k\right|^4
 \leq C''\ell^2.
\tag{15}
\]

Equation (9) also gives
\[
 \sup_\theta\mathbb E_\theta
 \left\lVert N_n^{-1}\sum_{k=0}^{N_n-1}U_k\right\rVert_2^2
 \leq C N_n^{-1}.
\tag{16}
\]
Therefore \(\widehat\xi_n\to\xi\) in probability uniformly on the margin class.  The same calculation with law-dependent constants gives convergence for each fixed law in \(\mathcal M_{\Delta,p}^{\mathrm{DAG},+}\).  Since \(\mathcal C\) in \(\text{def:centering-map}\) is a quadratic polynomial and hence continuous, \(\widehat\eta_n=\mathcal C(\widehat\xi_n)\to\mathcal C(\xi)=\eta=(m,\Gamma_1,\Gamma_2)\), proving the two asserted consistency claims.

We now prove the uniform Gaussian approximation.  For a unit vector \(a\), put \(Z_k=a^{\mathsf T}U_k\), choose \(\ell_n=\lfloor N_n^{1/2}\rfloor\) and \(g_n=\lceil c_0\log N_n\rceil\) with \(c_0>(2c_{\mathrm{mix}}\Delta)^{-1}\), and partition the observations into consecutive large blocks of length \(\ell_n\), separating gaps of length \(g_n\), plus a terminal remainder.  There are \(v_n\leq N_n/\ell_n\) large blocks.  The discarded set has size \(O(N_ng_n/\ell_n+\ell_n)\), so (9) makes its normalized sum converge to zero in \(L^2\), uniformly.  Iterating the defining strong-mixing event inequality, first for simple functions and then by bounded approximation, shows that replacing the large-block characteristic functions by those of independent block copies changes their product by at most
\[
 4v_n\alpha_{Y,\theta}(g_n)
 \leq C N_n^{1/2}e^{-c_{\mathrm{mix}}\Delta g_n}=o(1)
\tag{17}
\]
uniformly.

Let \(B_{b,n}\) denote the independent copies of the centered large-block sums.  From (15) and (9),
\[
 \sum_{b=1}^{v_n}\frac{\mathbb E|B_{b,n}|^4}{N_n^2}
 \leq C\frac{v_n\ell_n^2}{N_n^2}
 =O(\ell_n/N_n)=o(1),
\tag{18}
\]
and Cauchy--Schwarz gives
\[
 \sum_{b=1}^{v_n}\frac{\mathbb E|B_{b,n}|^3}{N_n^{3/2}}
 \leq C\frac{v_n\ell_n^{3/2}}{N_n^{3/2}}
 =O\{(\ell_n/N_n)^{1/2}\}=o(1).
\tag{19}
\]
For fixed \(u\in\mathbb R\), the elementary bound \(|e^{ix}-1-ix+x^2/2|\leq |x|^3/6\), (19), and \(\max_b\operatorname{Var}(B_{b,n})/N_n\leq C\ell_n/N_n=o(1)\) show, after taking logarithms of the product, that
\[
 \prod_{b=1}^{v_n}\mathbb E\exp(iuB_{b,n}/\sqrt{N_n})
 =\exp\{-u^2\sigma_{a,n}^2/2\}+o(1),
 \quad
 \sigma_{a,n}^2=N_n^{-1}\sum_{b=1}^{v_n}\operatorname{Var}(B_{b,n}),
\tag{20}
\]
uniformly in \(\theta\).  If \(\gamma_a(h)=\operatorname{Cov}(Z_0,Z_h)\), stationarity gives
\[
 \ell_n^{-1}\operatorname{Var}\left(\sum_{k=1}^{\ell_n}Z_k\right)
 =\sum_{|h|<\ell_n}(1-|h|/\ell_n)\gamma_a(h).
\tag{21}
\]
The geometric envelope (8), together with \(v_n\ell_n/N_n\to1\), makes (21) converge uniformly to \(a^{\mathsf T}\Omega(\theta)a\).  Combining (17), (20), (21), and the negligible discarded observations proves the scalar characteristic-function approximation uniformly for every fixed \(a\).

To obtain the stated vector bounded-Lipschitz conclusion without imposing compactness on the parameter set, argue by contradiction.  If it failed, choose \(\theta_n\) along which the discrepancy is bounded away from zero.  The uniformly bounded positive-semidefinite matrices \(\Omega(\theta_n)\) have a convergent subsequence, say to \(\Omega_*\).  The preceding scalar calculation applied to \(a=u/\lVert u\rVert_2\) for each \(u\in\mathbb R^{d_0}\) gives pointwise convergence of the vector characteristic functions to \(\exp(-u^{\mathsf T}\Omega_*u/2)\).  Equation (16) gives uniform tightness.  On a sufficiently large ball, the unit bounded-Lipschitz class has a finite uniform sup-norm net; continuous compactly supported approximations to the finitely many net members are approximated uniformly by functions with integrable Fourier transforms, for which characteristic-function convergence implies convergence of expectations.  Outside the ball, (16) controls both the empirical laws and the Gaussian laws.  Hence the bounded-Lipschitz discrepancy along the subsequence tends to zero, a contradiction.  Thus
\[
 \sup_\theta d_{\mathrm{BL}}\left(
 \mathcal L_\theta\{\sqrt n(\widehat\xi_n-\xi(\theta))\},
 \mathcal N\{0,\Omega(\theta)\}\right)\longrightarrow0,
\tag{22}
\]
where \(N_n/n\to1\) justifies replacing \(N_n\) by \(n\).

For \(h=\widehat\xi_n-\xi\), the exact quadratic expansion of \(\mathcal C\) is
\[
 \mathcal C(\xi+h)-\mathcal C(\xi)
 =D\mathcal C_\xi h+R(h),
 \qquad \lVert R(h)\rVert_2\leq C\lVert h\rVert_2^2.
\tag{23}
\]
The population first and second raw moments, and hence \(D\mathcal C_\xi\), are uniformly bounded by (1).  Equation (16) implies \(\sqrt nR(h)=o_{P_\theta}(1)\) uniformly.  Applying the uniformly bounded linear maps \(D\mathcal C_{\xi(\theta)}\) to (22), and then using (23), proves the asserted uniform Gaussian approximation for \(\sqrt n(\widehat\eta_n-\eta)\) with covariance \(D\mathcal C_\xi\Omega D\mathcal C_\xi^{\mathsf T}\).

It remains to prove the HAC claim.  For coordinates \(a,b,c,d\) and arbitrary integers \(h,h'\geq0\), the centered fourth-cumulant identity gives
\[
 \begin{aligned}
 &\operatorname{Cov}(U_0^aU_h^b,U_t^cU_{t+h'}^d)\\
 &\quad=\operatorname{cum}(U_0^a,U_h^b,U_t^c,U_{t+h'}^d)
 +\operatorname{Cov}(U_0^a,U_t^c)\operatorname{Cov}(U_h^b,U_{t+h'}^d)\\
 &\qquad\quad
 +\operatorname{Cov}(U_0^a,U_{t+h'}^d)\operatorname{Cov}(U_h^b,U_t^c).
 \end{aligned}
\tag{24}
\]
Order the multiset \(\{0,h,t,t+h'\}\), and let \(G_t\) be its largest adjacent gap.  Equation (13) bounds the cumulant in (24) by \(C\rho^{G_t/2}\).  For every integer \(q\geq0\), if \(G_t\leq q\), the diameter of these four ordered points is at most \(3q\); because one point is zero, this forces \(|t|\leq3q\).  Thus at most \(6q+1\) integers \(t\) have \(G_t\leq q\), uniformly in \(h,h'\), and consequently
\[
 \sup_{h,h'\geq0}\sum_{t\in\mathbb Z}\rho^{G_t/2}
 \leq 1+\sum_{q=1}^\infty(6q+1)\rho^{(q-1)/2}<\infty.
\tag{25}
\]
For each of the other two terms in (24), (8) gives a convolution of one absolutely summable covariance sequence with a uniformly bounded one.  For example,
\[
 \sum_{t\in\mathbb Z}
 |\operatorname{Cov}(U_0^a,U_t^c)
   \operatorname{Cov}(U_h^b,U_{t+h'}^d)|
 \leq
 \lVert\gamma_{ac}\rVert_1\lVert\gamma_{bd}\rVert_\infty\leq C.
\tag{26}
\]
Equations (24)--(26) prove the required translation-uniform product bound
\[
 \sup_{h,h'\geq0}\sum_{t\in\mathbb Z}
 |\operatorname{Cov}(U_0^aU_h^b,U_t^cU_{t+h'}^d)|<\infty.
\tag{27}
\]

Define the true-mean lag average
\[
 \widetilde\Xi_h=N_n^{-1}\sum_{k=0}^{N_n-h-1}U_kU_{k+h}^{\mathsf T}.
\]
For all sufficiently large \(n\), \(b_n<N_n\).  Expanding the two finite averages and setting \(t=\ell-k\), stationarity and (27) yield, uniformly for \(0\leq h,h'\leq b_n\),
\[
 \begin{aligned}
 &\left|\operatorname{Cov}\{(\widetilde\Xi_h)_{ab},
                         (\widetilde\Xi_{h'})_{cd}\}\right|\\
 &\quad\leq N_n^{-2}\sum_{k=0}^{N_n-h-1}
       \sum_{\ell=0}^{N_n-h'-1}
       |\operatorname{Cov}(U_k^aU_{k+h}^b,U_\ell^cU_{\ell+h'}^d)|\\
 &\quad\leq N_n^{-2}\sum_{k=0}^{N_n-h-1}
       \sum_{t\in\mathbb Z}
       |\operatorname{Cov}(U_0^aU_h^b,U_t^cU_{t+h'}^d)|
 \leq C/N_n.
 \end{aligned}
\tag{28}
\]
There are only \(d_0^2\) matrix coordinates, the Bartlett weights lie in \([0,1]\), and symmetrization changes only a fixed factor.  Summing (28) over the \((b_n+1)^2\) lag pairs gives
\[
 \sup_\theta\mathbb E_\theta\left\lVert
 \widetilde\Xi_0-\mathbb E\widetilde\Xi_0
 +\sum_{h=1}^{b_n}w(h/b_n)
 \{\widetilde\Xi_h+\widetilde\Xi_h^{\mathsf T}
   -\mathbb E(\widetilde\Xi_h+\widetilde\Xi_h^{\mathsf T})\}
 \right\rVert_{\mathrm F}^2
 \leq Cb_n^2/N_n=o(1).
\tag{29}
\]

Let \(\Xi_h=\operatorname{Cov}(Y_0,Y_h)\).  Directly from stationarity,
\[
 \mathbb E\widetilde\Xi_h=(1-h/N_n)\Xi_h.
\tag{30}
\]
The geometric envelope (8) is a uniform summable envelope for \(\Xi_h\).  Since \(w(h/b_n)\to1\) for every fixed \(h\), equations (9) and (30), split at a fixed lag and then with that lag sent to infinity, show uniformly that the expectation of the symmetrized expression in (29) converges to
\[
 \Xi_0+\sum_{h=1}^\infty(\Xi_h+\Xi_h^{\mathsf T})
 =\sum_{h\in\mathbb Z}\operatorname{Cov}(Y_0,Y_h)=\Omega(\theta).
\tag{31}
\]

Finally put \(\overline U_n=N_n^{-1}\sum_{k=0}^{N_n-1}U_k=\widehat\xi_n-\xi\), and define
\[
 A_{h,n}=N_n^{-1}\sum_{k=0}^{N_n-h-1}U_k,
 \qquad
 B_{h,n}=N_n^{-1}\sum_{k=0}^{N_n-h-1}U_{k+h}.
\]
The empirical lag matrix in \(\text{def:hac-estimator}\) satisfies the exact expansion
\[
 \widehat\Xi_h-\widetilde\Xi_h
 =-A_{h,n}\overline U_n^{\mathsf T}
  -\overline U_nB_{h,n}^{\mathsf T}
  +(1-h/N_n)\overline U_n\overline U_n^{\mathsf T}.
\tag{32}
\]
By (16), \(\lVert\overline U_n\rVert_2=O_{P_\theta}(N_n^{-1/2})\) uniformly.  The first moment bound from (2) and Markov's inequality give
\[
 \sum_{h=0}^{b_n}\{\lVert A_{h,n}\rVert_2+\lVert B_{h,n}\rVert_2\}=O_{P_\theta}(b_n)
\tag{33}
\]
uniformly.  Therefore (32), after Bartlett weighting and symmetrization, changes the complete lag-window sum by
\[
 O_{P_\theta}\{b_nN_n^{-1/2}+b_nN_n^{-1}\}=o_{P_\theta}(1)
\tag{34}
\]
uniformly, using \(\text{ass:hac-bandwidth}\).  The empty-sum convention covers the finitely many samples before \(b_n<N_n\).  Combining (29), (31), and (34), then using \(\lVert M\rVert_{\mathrm{op}}\leq\lVert M\rVert_{\mathrm F}\), proves for every \(\varepsilon>0\),
\[
 \sup_{\theta:P_\theta\in\mathcal M_{\Delta,p}(\kappa)}
 P_\theta\{\lVert\widehat\Omega_n-\Omega(\theta)\rVert_{\mathrm{op}}>\varepsilon\}
 \longrightarrow0.
\]